\section{Additional Results and Mechanism Evaluation}
\label{app:additional_results}

\subsection{Task-Level Robustness to Misleading Experience}
\label{app:fullresults}

Table~\ref{tab:app_distraction_tasks} resolves Figure~\ref{fig:experience_distraction} into its two tasks.
Similar-event replacement reduces full \method{} by 10.8 points on Put Back Block and 10.6 on Press Button, while the no-comparison drops are 30.4 and 30.6 points.
The resulting degradation advantages, 19.6 and 20.0 points, show the same robustness pattern across both tasks.
Relative to the recollection-only reference, \method{} remains ahead by 15.8 and 17.4 points, while the no-comparison variant falls behind by 8.0 and 5.6 points.

\begin{table}[H]
\centering\small
\caption{\textbf{Task-level experience-distraction results (\%).} Mean $\pm$ sample SD across five evaluation runs, with 100 rollouts per task/run. The recollection-only row is the normal-condition reference. Bold/underlined means are best/second-best within each retrieval condition.}
\label{tab:app_distraction_tasks}
\begin{tabular}{@{}llcc@{}}
\toprule
Method & Retrieved events & Put Back Block $\uparrow$ & Press Button $\uparrow$ \\
\midrule
\method{} & Normal & $\mathbf{100.0}\pm0.0$ & $\mathbf{92.0}\pm1.6$ \\
 & Random distractors & $\mathbf{94.6}\pm3.4$ & $\mathbf{86.6}\pm4.8$ \\
 & Similar distractors & $\mathbf{89.2}\pm4.5$ & $\mathbf{81.4}\pm5.4$ \\
\midrule
Without explicit comparison & Normal & $\underline{95.8}\pm2.2$ & $\underline{89.0}\pm2.5$ \\
 & Random distractors & $\underline{82.6}\pm6.5$ & $\underline{75.8}\pm6.7$ \\
 & Similar distractors & $\underline{65.4}\pm6.0$ & $\underline{58.4}\pm10.6$ \\
\midrule
Recollection only & Reference & $73.4\pm4.0$ & $64.0\pm8.7$ \\
\bottomrule
\end{tabular}
\end{table}

The two tasks require recollection of an original destination and repeated-interaction progress, respectively.
Across both demands, explicit comparison helps preserve the benefit of recalled experience.
We next examine the candidate-level evidence underlying this task-level comparison.

\subsection{Candidate Assessment and Selective Admission}
\label{app:mechanism_probes}

Our evidence chain follows the decisions made by \method{}: predict a proposal's effect, assess its task relevance, and regulate its influence on generation.
We use shared branch execution to evaluate effect accuracy and task-appropriate choice, then check joint withdrawal in the generation interface.


\subsubsection{Effect Accuracy and Task-Appropriate Selection}
\label{app:effect_probes}

\paragraph{Shared execution and independent outcomes.}
Our branch protocol fixes the checkpoint, source episodes, stages, and queries before scoring.
For each adapted proposal $\mu_i$, execution starts from restored physical, controller, task-progress, random-generator, and observation/history states.
Executing $H=32$ commands without replanning yields the target
\begin{equation}
 E^{\mathrm{obs}}_{t,i}=\mathcal P_{\mathrm{eff}}\!\left(
 \Phi(o_t),\Phi(o^{\mu_i}_{t+H})-\Phi(o_t)\right).
 \label{eq:effect_probe_target}
\end{equation}
Effect metrics use branches with complete $H$-step endpoints; applicability uses independently resolvable task outcomes, including early terminal outcomes when the task criterion is met.
We compare the three effect references on the same valid endpoints and labels for each metric.
We assign $a_{t,i}\in\{0,1\}$ by a predefined next-interaction success criterion or blinded outcome review, retaining unresolved labels as unknown.
Outcome review supplies independent labels for both consequence prediction and candidate admission.

\paragraph{Effect accuracy.}
We compare $F_{t,i}\in\{E_i,r_t,\widehat E_i\}$ with $E^{\mathrm{obs}}_{t,i}$ on shared valid endpoints using the feature-width-normalized MSE $d_E$ in Eq.~\eqref{eq:distances}.
The three references distinguish a stored historical effect, a candidate-independent requirement, and an action-conditioned prediction in the current world context.
We average candidates within query, then queries within source episode, episodes within task, and tasks with equal weights.

\paragraph{Ordering and applicable choice.}
Let $s^F_{t,i}=\kappa(F_{t,i},r_t)$ and $s^{\mathrm{obs}}_{t,i}=\kappa(E^{\mathrm{obs}}_{t,i},r_t)$.
The comparable set $\mathcal P_t$ contains candidate pairs with observed-score differences exceeding $\delta_{\mathrm{obs}}$.
With $\operatorname{sgn}_\delta(x)=0$ for $|x|\leq\delta$ and $\operatorname{sign}(x)$ otherwise, ordering agreement is
\begin{equation}
 h^F_{t,ij}=\begin{cases}
 1/2,&\operatorname{sgn}_{\delta_{\mathrm{pred}}}(s^F_{t,i}-s^F_{t,j})=0,\\
 1,&\operatorname{sgn}_{\delta_{\mathrm{pred}}}(s^F_{t,i}-s^F_{t,j})
       =\operatorname{sign}(s^{\mathrm{obs}}_{t,i}-s^{\mathrm{obs}}_{t,j}),\\
 0,&\text{otherwise}.
 \end{cases}
 \label{eq:tie_ordering}
\end{equation}
We average pairs within query, then queries, episodes, and tasks as above; tie rate is the fraction of comparable pairs with tied predicted scores.
Queries without comparable pairs are excluded from ordering agreement.
Development repeatability checks determine all tolerances before held-out evaluation.

For the maximum-score tie set $\mathcal T_t^F$, applicable choice is
$|\mathcal T_t^F|^{-1}\sum_{i\in\mathcal T_t^F}a_{t,i}$ when all its labels are known, matching uniform tie breaking.
On comparable pairs, query-only scores yield $1/2$ ordering agreement and 100\% ties analytically; applicable choice equals the candidate-pool applicability rate, which requires outcome labels.
These metrics connect requirement-aligned ordering to independently verified task progress.

\begingroup
\raggedbottom

\subsubsection{Candidate-Specific Effects and Full-Selector Quality}
\label{app:rebuttal_selector}

This study extends Table~\ref{tab:effect_probes} and Appendix~\ref{app:effect_probes} from effect accuracy to decision quality.
The central question is whether $\widehat E_i$ distinguishes actions within the same candidate pool and improves the deployed selector, rather than only predicting a query's average future.
The expanded evaluation uses a separate cohort and complements, rather than replaces, the main-text diagnostic.

\paragraph{Shared queries, branches, and outcome labels.}
The cohort comprises 20 held-out source episodes per task on Put Back Block (PBB), Press Button (PB), and Battery Try (BT), with ten prespecified replans per episode and the deployed $K=32$ candidate slots at each replan.
Episodes used for fitting or hyperparameter/checkpoint selection are excluded.
The query schedule is fixed before inspecting predictions: five replans come from approach/alignment stages and five from contact/completion stages of the relevant interaction, using simulator task predicates rather than model scores to identify stages.
Episodes with fewer available queries are retained with their actual counts; they are not replaced according to outcomes.
At each query, we cache the requirement, adapted proposals, rank scores, and cursor state once.
We execute each proposal for $H=32$ commands from the same restored simulator, controller, task-progress, RNG, and factual-history state, following Appendix~\ref{app:effect_probes}; no Action DiT refinement is applied in this branch diagnostic.

Applicability is a task-progress label, not a threshold on $\kappa_i$.
For PBB, a branch must advance the correct grasp/return interaction toward the episode's original pad, without placing the block on a wrong pad.
For PB, it must advance approach to the required button or register a permitted press without exceeding the remaining count.
For BT, it must advance alignment/insertion of the battery currently required by task feedback, without repeating a feedback-rejected attempt.
Before evaluation, each sampled stage is assigned one explicit predicate: reaching the task evaluator's contact/alignment region, acquiring the correct object, or completing the required placement/press/insertion.
The mapping and geometric tolerances are frozen on development episodes; a branch is positive only if that predicate is reached without an invalid transition.
No progress or an incorrect interaction gives $a_{t,i}=0$; an outcome that cannot be resolved from the simulator state is unknown, not automatically negative.
Labels use privileged state only for evaluation, never as policy or selector inputs.

\begin{table}[H]
\centering\small
\caption{\textbf{Branch-cohort accounting.}
Full endpoints, early task termination, and missing endpoints partition attempted branches.
Applicability counts are separate: an early terminal outcome can still have a known label.
Coverage and choice accuracy use the same fully labelled query cohort.}
\label{tab:rebuttal_coverage}
\begin{tabular}{@{}lrrr@{}}
\toprule
Accounting item & PBB & PB & BT \\
\midrule
Source episodes / queries & 20 / 200 & 20 / 200 & 20 / 200 \\
Attempted candidate branches & 6,400 & 6,400 & 6,400 \\
Complete $H$-step endpoints & 6,304 & 6,272 & 6,208 \\
Early task termination & 96 & 128 & 192 \\
Missing endpoints / unknown labels & 0 / 0 & 0 / 0 & 0 / 0 \\
Applicable / inapplicable labels & 1,152 / 5,248 & 1,024 / 5,376 & 896 / 5,504 \\
Common fully labelled queries & 200 & 200 & 200 \\
Queries with at least two full endpoints & 200 & 200 & 200 \\
Queries with comparable effect pairs & 196 & 192 & 186 \\
Queries with an applicable candidate & 100 & 80 & 70 \\
Applicable-pool coverage (\%) & 50.0 & 40.0 & 35.0 \\
\bottomrule
\end{tabular}
\end{table}

\paragraph{Candidate-specific effect error.}
For the shared complete-endpoint candidate set $\mathcal V_t$ at query $t$, center each reference $F_{t,i}\in\{E_i,r_t,\widehat E_i\}$ and its observed target within that same set:
\begin{equation}
 \widetilde F_{t,i}=F_{t,i}-\frac{1}{|\mathcal V_t|}\sum_{j\in\mathcal V_t}F_{t,j},\quad
 \widetilde E^{\mathrm{obs}}_{t,i}=E^{\mathrm{obs}}_{t,i}
 -\frac{1}{|\mathcal V_t|}\sum_{j\in\mathcal V_t}E^{\mathrm{obs}}_{t,j}.
 \label{eq:rebuttal_centering}
\end{equation}
The centered error is $|\mathcal V_t|^{-1}\sum_{i\in\mathcal V_t}d_E(\widetilde F_{t,i},\widetilde E^{\mathrm{obs}}_{t,i})$, with $d_E$ from Eq.~\eqref{eq:distances}; queries with fewer than two complete endpoints are excluded from this metric and counted separately.
Because the centered query-only prediction is zero, this diagnostic tests candidate differences rather than the common query component.
Ordinary MSE and ordering agreement retain the definitions in Appendix~\ref{app:effect_probes}.
Development-branch repeatability checks fix ordering tolerances before held-out evaluation.

\begin{table}[H]
\centering\small
\caption{\textbf{Effect prediction within shared queries.}
Centering removes each query's mean effect. Macro averages weight tasks equally; query-only ordering is exactly $1/2$ on comparable pairs by definition, not an empirical gain.}
\label{tab:rebuttal_effects}
\begin{tabular}{@{}llccc@{}}
\toprule
Task & Effect reference & MSE $\downarrow$ & Centered MSE $\downarrow$ & Ordering $\uparrow$ \\
\midrule
PBB & Historical $E_i$ & 0.050 & 0.036 & 0.430 \\
 & Query-only $r_t$ & 0.039 & 0.035 & 0.500 \\
 & Predicted $\widehat E_i$ & 0.028 & 0.021 & 0.670 \\
PB & Historical $E_i$ & 0.058 & 0.039 & 0.410 \\
 & Query-only $r_t$ & 0.043 & 0.039 & 0.500 \\
 & Predicted $\widehat E_i$ & 0.033 & 0.024 & 0.640 \\
BT & Historical $E_i$ & 0.063 & 0.044 & 0.380 \\
 & Query-only $r_t$ & 0.048 & 0.043 & 0.500 \\
 & Predicted $\widehat E_i$ & 0.037 & 0.029 & 0.600 \\
\midrule
Macro & Historical $E_i$ & 0.057 & 0.040 & 0.407 \\
 & Query-only $r_t$ & 0.043 & 0.039 & 0.500 \\
 & Predicted $\widehat E_i$ & 0.033 & 0.025 & 0.637 \\
\bottomrule
\end{tabular}
\end{table}

\paragraph{Separate retrieval coverage from selection quality.}
For selection, we use one common set $\mathcal Q$ of queries whose entire valid candidate pool has known applicability labels, including known early-terminal outcomes.
Candidates are not removed because their predicted scores are poor or their $H$-step endpoints are unavailable.
Queries with unknown labels remain in the accounting; we additionally report all-query lower/upper bounds by assigning unknown outcomes 0/1.
Let $p^s_{t,i}$ be the selection probability for rule $s$, uniform over its maximum-score ties (absolute score tolerance $10^{-8}$), and let $\mathcal Q^+=\{t\in\mathcal Q:\max_i a_{t,i}=1\}$.
Define per-query choice accuracy $c_t^s=\sum_i p^s_{t,i}a_{t,i}$, pool coverage $\Pr(t\in\mathcal Q^+)$, all-query accuracy $\mathbb E[c_t^s]$, and covered-pool accuracy $\mathbb E[c_t^s\mid t\in\mathcal Q^+]$.
We use equal episode weight within each task, assigning equal weight to queries within episode; the conditional metric renormalizes those same query weights on $\mathcal Q^+$.
Within each task, all-query accuracy equals coverage times conditional accuracy, separating retrieval coverage from candidate selection quality.

The five rules are uniform random selection; historical compatibility $\kappa(E_i,r_t)$; predicted compatibility $\kappa(\widehat E_i,r_t)$; contextual selection $\ell_i^{\mathrm{rank}}+b_{t,i}$; and the full selector in Eq.~\eqref{eq:selection}.
They share the actual candidate pool, adapted actions, $r_t$, rank scores, and continuity state; only the scoring rule changes, and no counterfactual selection updates the cached cursor.
Removing $\kappa_i$ here is a \emph{fixed-checkpoint selector diagnostic}, not the separately configured no-comparison policy in Table~\ref{tab:ablation}.

\begin{table}[H]
\centering\footnotesize
\setlength{\tabcolsep}{4pt}
\caption{\textbf{Independent applicability of selected actions (\%).}
Each cell is all-query / covered-pool accuracy. All rules use the same labelled pools; the macro conditional column averages the task-conditional rates, not a pooled ratio.}
\label{tab:rebuttal_selection}
\begin{tabular}{@{}lcccc@{}}
\toprule
Selection rule & PBB & PB & BT & Macro \\
\midrule
Uniform random & 18.0 / 36.0 & 16.0 / 40.0 & 14.0 / 40.0 & 16.0 / 38.7 \\
Historical compatibility & 20.0 / 40.0 & 17.0 / 42.5 & 15.0 / 42.9 & 17.3 / 41.8 \\
Predicted compatibility & 23.0 / 46.0 & 20.0 / 50.0 & 17.0 / 48.6 & 20.0 / 48.2 \\
$\ell^{\mathrm{rank}}+b$ & 27.0 / 54.0 & 23.0 / 57.5 & 19.0 / 54.3 & 23.0 / 55.3 \\
Full selector & 31.0 / 62.0 & 28.0 / 70.0 & 23.0 / 65.7 & 27.3 / 65.9 \\
\bottomrule
\end{tabular}
\end{table}

\paragraph{Paired uncertainty and interpretation.}
We resample source episodes within each task 10,000 times, keeping all queries, branches, and scoring rules of each sampled episode together.
We report percentile 95\% intervals for paired differences and the equal-task mean, treating source episodes rather than candidate pairs as independent units.
Table~\ref{tab:rebuttal_effect_deltas} reports the principal contrasts, with task-wise directions visible in the preceding tables.

\begin{table}[H]
\centering\small
\caption{\textbf{Paired mechanism contrasts.}
Brackets give episode-clustered bootstrap 95\% intervals.
Centered-MSE contrasts are predicted minus reference; choice contrasts are percentage points.}
\label{tab:rebuttal_effect_deltas}
\begin{tabularx}{\linewidth}{@{}Xr@{}}
\toprule
Contrast & Estimate [95\% interval] \\
\midrule
Centered MSE: predicted minus query-only & $-0.014\ [-0.019,-0.009]$ \\
Centered MSE: predicted minus historical & $-0.015\ [-0.020,-0.009]$ \\
All-query choice: predicted minus historical compatibility & $+2.7\ [0.9,4.4]$ \\
All-query choice: full selector minus $\ell^{\mathrm{rank}}+b$ & $+4.3\ [2.0,6.7]$ \\
\bottomrule
\end{tabularx}
\end{table}

\paragraph{Candidate discrimination and decision quality.}
The predicted-versus-query-only centered-error difference is -0.014, while the full selector changes all-query choice accuracy by +4.3 percentage points relative to $\ell^{\mathrm{rank}}+b$ (Table~\ref{tab:rebuttal_effect_deltas}).
These contrasts distinguish two claims: sensitivity to candidate-specific consequences, and a benefit of that sensitivity for choosing actions.
A lower centered error together with a positive paired choice contrast supports both, whereas lower effect error alone did not establish better selection.
All-query accuracy (27.3\%) is interpreted jointly with pool coverage (41.7\%) and covered-pool accuracy (65.9\%), separating unavailable experience from selection errors.
These are pre-refinement diagnostics, not rollout success rates or evidence that every rejected proposal is useless after refinement.

\subsubsection{Adaptive Utility Beyond Deterministic Admission Rules}
\label{app:rebuttal_gate}

This study complements Table~\ref{tab:ablation} and Figure~\ref{fig:experience_distraction} by asking whether learned utility regulates recalled experience in practice, beyond the structural zero-gate property tested in Appendix~\ref{app:gate_probes}.

\paragraph{Record quality and gate values before refinement.}
We evaluate normal retrieval and the same random/similar complete-event replacements as Appendix~\ref{app:distraction_protocol} on PBB and PB.
For each condition, we prespecify ten replans from each of 20 held-out episodes per task and log the full learned-gate policy's selected proposal, raw $\alpha_t$, effective $g_t$, and all rule flags before executing the refined action.
We restore each sampled state to execute the selected adapted proposal and label its applicability using the independent task predicates above.
No training-negative/null flag is artificially enabled during evaluation.
Empty pools have $g_t=0$ but no selected-candidate applicability label and are recorded separately, without assigning an arbitrary $\alpha_t$ or negative label.
Unknown branch labels likewise remain in coverage counts.
Empty-pool / unknown-label counts are 0/0 for normal, 0 for random, and 0 for similar retrieval.

To separate learning from rules, define $v_t^{\mathrm{det}}=1$ when the selected proposal passes availability/provenance, temporal eligibility, and the hard repetition cutoff, \emph{before} the learned threshold $\alpha_{\min}$.
Raw-utility summaries and learned rejection rates are conditioned on $v_t^{\mathrm{det}}=1$; effective-gate summaries include deterministic vetoes for every labelled selected proposal.
This prevents a lower mean gate caused by hard rules from being attributed to learned discrimination.

\begin{table}[H]
\centering\footnotesize
\setlength{\tabcolsep}{3pt}
\caption{\textbf{Gate response to selected-proposal quality.}
Medians [Q1, Q3] replace a distribution plot. $N/N_{\mathrm{det}}$ gives labelled/eligible counts.
The final column is $\Pr(\alpha_t<\alpha_{\min}\mid v_t^{\mathrm{det}}=1,a)$ in percent, with the deployed threshold unchanged.
Empty pools and unknown labels are reported separately.}
\label{tab:rebuttal_gate_quality}
\begin{tabular}{@{}llrccc@{}}
\toprule
Retrieval & Applicability & $N/N_{\mathrm{det}}$ & $\alpha_t\mid v^{\mathrm{det}}=1$ & $g_t$ & Rejected (\%) \\
\midrule
Normal & Applicable & 260 / 252 & .72 [.56, .84] & .55 [.38, .68] & 4.0 \\
 & Inapplicable & 140 / 116 & .23 [.10, .39] & .04 [.00, .17] & 34.5 \\
Random & Applicable & 220 / 208 & .70 [.54, .82] & .53 [.35, .66] & 4.8 \\
 & Inapplicable & 180 / 140 & .19 [.08, .33] & .00 [.00, .13] & 46.4 \\
Similar & Applicable & 200 / 188 & .68 [.52, .81] & .51 [.32, .65] & 5.3 \\
 & Inapplicable & 200 / 148 & .14 [.05, .28] & .00 [.00, .08] & 54.1 \\
\bottomrule
\end{tabular}
\end{table}

\begin{table}[H]
\centering\small
\setlength{\tabcolsep}{4pt}
\caption{\textbf{Why candidate influence becomes zero (\%).}
First-veto attribution uses the displayed left-to-right order, so causes are mutually exclusive.
All rates use the full logged-query cohort per condition; the complete overlapping flag vector is retained in the logs.
Repetition's continuous attenuation is reflected in $g_t$, not counted as a hard veto.}
\label{tab:rebuttal_veto}
\begin{tabular}{@{}lccccc@{}}
\toprule
Retrieval & Other eligibility & Continuity & Repetition & $\alpha<\alpha_{\min}$ & Active $g>0$ \\
\midrule
Normal & 1.0 & 4.0 & 3.0 & 12.5 & 79.5 \\
Random & 2.0 & 6.0 & 5.0 & 18.8 & 68.2 \\
Similar & 3.0 & 7.0 & 6.0 & 22.5 & 61.5 \\
\bottomrule
\end{tabular}
\end{table}

The tables condition on the \emph{selected} proposal: a distractor-rich pool does not require a lower gate if selection still finds useful experience.
Moreover, $\alpha_t$ is trained as imitation utility, not a calibrated probability that a raw proposal is applicable.

\paragraph{A matched constant-utility control.}
We train a control that substitutes a development-selected constant $\alpha_0$ for the learned utility in \emph{both} source and conditioning, while retaining identical deterministic eligibility, continuity, hard/soft repetition, retrieval, selector architecture, losses, and training budget.
Candidate availability and eligibility checks, including the comparison to $\alpha_{\min}$, operate on the substituted value.
We choose $\alpha_0$ from $\{0.25,0.50,0.75\}$ using short, equally budgeted development-only runs and the equal-task/equal-condition mean SR; ties choose the smaller value.
This choice is frozen before full-budget training and held-out evaluation; the selected constant is $\alpha_0=\text{0.50}$.
The control retains the utility head and its auxiliary supervision but disconnects its output from action generation, isolating use of adaptive utility rather than removal of a loss or module.
Initialization, data ordering, optimizer schedules, NFE, action horizon/prefix, corruption schedules, and checkpoint-selection rules are matched.
Each variant is trained under its own deployed gate rule with three paired training seeds (11, 22, 33).
The retrieval/selector procedure is unchanged, but its fitted weights and visited states may differ across trained policies; only the fixed-query diagnostic above holds numerical scores and candidates identical.

Each trained checkpoint is evaluated on 100 paired resets per task/condition, with identical donor-sampling and policy-noise schedules at reset and the existing task step limits.
This gives 3,600 evaluation episodes across two variants, three seeds, two tasks, and three conditions, excluding development selection.
Subsequent trajectories are allowed to diverge; no shared on-policy candidate sequence is assumed.
We report mean $\pm$ sample SD across \emph{training seeds}, after averaging reset outcomes within seed.
This control differs from the no-comparison model, and its training-seed variability differs from the five evaluation-run variability in Figure~\ref{fig:experience_distraction}.

\begin{table}[H]
\centering\small
\setlength{\tabcolsep}{4pt}
\caption{\textbf{Learned versus constant utility (\%).}
Mean $\pm$ sample SD uses three training seeds; macro scores average both tasks within each seed before computing SD.
Drop is normal minus similar-distraction SR, computed within seed.}
\label{tab:rebuttal_constant_gate}
\begin{tabular}{@{}llcccc@{}}
\toprule
Task & Utility & Normal & Random & Similar & Similar drop $\downarrow$ \\
\midrule
PBB & Learned & $99.3\pm0.6$ & $94.0\pm2.0$ & $89.0\pm1.0$ & $10.3\pm0.6$ \\
 & Constant & $99.0\pm1.0$ & $89.7\pm1.5$ & $78.7\pm2.1$ & $20.3\pm1.5$ \\
PB & Learned & $92.0\pm1.0$ & $86.3\pm1.5$ & $81.3\pm1.2$ & $10.7\pm1.5$ \\
 & Constant & $91.3\pm1.5$ & $80.7\pm2.1$ & $69.7\pm2.1$ & $21.7\pm1.2$ \\
\midrule
Macro & Learned & $95.7\pm0.8$ & $90.2\pm1.4$ & $85.2\pm1.0$ & $10.5\pm1.0$ \\
 & Constant & $95.2\pm1.0$ & $85.2\pm0.6$ & $74.2\pm0.6$ & $21.0\pm0.5$ \\
\bottomrule
\end{tabular}
\end{table}

For uncertainty, we retain paired reset outcomes within task and seed and report the three seed-level contrasts separately.
A descriptive $t$ interval on paired seed differences uses $\bar d\pm t_{0.975,2}s_d/\sqrt3$; with only three seeds, this summary depends on a normality assumption and has limited precision.
The similar-condition macro contrast is 11.0 points (three seed-level contrasts: 11.5,9.5,12.0; descriptive 95\% interval: [7.7,14.3]).
The reduction in the normal-to-similar drop is 10.5 points (95\% interval: [8.3,12.7]).

\paragraph{Adaptive strength versus structural withdrawal.}
The matched control distinguishes the value of adaptive strength from the shared rules that disable candidate influence.
Comparable normal-condition success together with a smaller distraction-induced drop supports a benefit beyond those unchanged rules.
Higher raw utility for applicable proposals among deterministically eligible cases additionally localize that response to learning rather than veto frequency.
The closed-loop result is essential: low raw-proposal applicability alone does not prove that refinement cannot make use of the recalled structure.
The structural withdrawal guarantee alone does not establish a benefit of learned adaptive strength; that claim depends on the matched closed-loop comparison.


\subsubsection{Shared Withdrawal Under Candidate Replacement}
\label{app:gate_probes}

We test the candidate-null boundary of Proposition~\ref{prop:withdrawal} on 50 fixed development prefixes from five Press Button and five Put Back Block episodes.
With current inputs, sampling randomness, and conditioning masks fixed, we replace retrieved candidates while forcing both the source and candidate-conditioning gates to zero.
The resulting $32\times14$ action outputs are identical in all 50 comparisons (maximum absolute difference 0; \texttt{v8-step300}, BF16, batch one, 20 inference steps).
This check verifies exact joint withdrawal of both candidate pathways under matched inputs.
\subsection{Real-World Tasks and Evaluation Protocol}
\label{app:realworld}

Banana-to-Box (B2B) and Duck-to-Box (D2B) probe repertoire reuse under changed layouts and poses (Section~\ref{sec:realworld}).
Return-to-Plate (R2P) and Same-Order-Retrieval (SOR) add episode-specific destination and order requirements for the same banana and rubber-duck objects, testing how recollection guides familiar interactions across task goals and scenes.
We evaluate 20 autonomous trials per method/task with shared resets and a fixed repertoire; recollection resets between trials.

\subsubsection{Platform, Data, and Shared Interface}
\label{app:realworld_hardware}

We use a single-arm AgileX PiPER-series manipulator with a parallel-jaw gripper.
Fast-WAM provides a same-family world--action reference, and $\pi_{0.5}$ provides a VLA reference under a shared physical action interface.
The comparison shares external (head) and wrist RGB views, proprioception, task instructions, training splits, reset schedules, and execution limits across methods.
Table~\ref{tab:realworld_hardware} records the data and adaptation configuration.

\paragraph{Recorded demonstrations and preprocessing.}
The current real-robot dataset contains 25 teleoperated demonstrations,
partitioned into 20 training and 5 development episodes.
The train/development counts are 4/1 for Banana-to-Box,
4/1 for Duck-to-Box, 8/2 for Return-to-Plate, and
4/1 for Same-Order-Retrieval.
The processed catalog contains 12,333 command-bearing frames:
9,976 in training and 2,357 in development.
The episode-level training/development split is distinct from the tabletop-variation evaluation in Figure~\ref{fig:realworld_c4_sequences}.
The repertoire contains 1,356 events from the 20 training episodes only.
Retrieval uses $K=32$ and a $768+4$-dimensional visual--task key.
We account for demonstrations and autonomous policy trials separately.

\begin{table}[H]
\centering\small
\caption{\textbf{Physical-data and adaptation configuration.}
Data timing describes the demonstration stream; deployment actions use the interface specified below.}
\label{tab:realworld_hardware}
\begin{tabularx}{\linewidth}{@{}p{0.25\linewidth}X@{}}
\toprule
Component & Recorded setting \\
\midrule
Platform & PiPER-series single arm with parallel-jaw gripper. \\
Observation & External/head and wrist RGB; declared $640\times480$ inputs, resized to $224\times224$ per view and joined into a $224\times448$ image. \\
Proprioception & Seven dimensions: TCP position, rotation vector, and gripper width. \\
Action target & Six absolute joint targets in radians and absolute gripper width in meters. \\
Temporal interface & Nominal data rate 20\,Hz; prediction horizon $H=32$. \\
Training / repertoire & 20 train / five development demonstrations; 1,356 training-only events; $K=32$. \\
Adaptation configuration & BF16; per-device batch one, accumulation one; learning rate $10^{-4}$ with cosine decay; weight decay 0.01; gradient clipping 1.0. \\
Randomness / budget & Training seed 42; corruption seed 3407; eight epochs configured. \\
\bottomrule
\end{tabularx}
\end{table}

\paragraph{Consistent state and action semantics.}
The joint-action conversion retains TCP-based proprioception and republishes joint targets as actions.
Training normalization statistics and the event bank are rebuilt for this action representation.
Although both vectors have seven dimensions, state and action retain their respective Cartesian and joint-control semantics.
For absolute-joint alignment, the control-space starts $c_t,c_i^0$ refer to joint coordinates; TCP proprioception $q_t$ remains a distinct input to the learned adapters.
The 20\,Hz value describes demonstration sampling; the joint targets are absolute commands rather than frame-to-frame differences.

The adaptation entries specify the training settings; the nominal data frequency and epoch budget retain their configuration-level meaning.



\begin{figure}[!t]
\vspace{-0.4cm}
\centering
\includegraphics[width=0.96\linewidth]{fig/physical_scenarios.pdf}
\caption{\textbf{Seven physical deployment sequences with retrieved cross-task events.}
Within each task pair, the upper row shows left-to-right execution keyframes from the external/head camera with wrist-view insets; the lower \emph{Past recall} row shows observations retrieved from indexed events in other tasks.
Original and Black identify the two tabletop appearances.
Each pair exposes a familiar interaction recalled under a new task or scene.}
\vspace{-0.4cm}
\label{fig:realworld_c4_sequences}
\end{figure}







\paragraph{Shared repertoire and traceable reuse.}
We build one repertoire from the 20 training episodes across all four tasks.
Inspection of its event indices confirms 1,356 training events and excludes the five development episodes.
Normalization uses the training partition.
Task identity contributes to retrieval similarity (Appendix~\ref{app:gate_details}); source episode and event identifiers preserve action provenance.

Each indexed event contains its source task, episode, action interval, and observation history.
The retrieved cross-task observations in Figure~\ref{fig:realworld_c4_sequences} make repertoire reuse visible, while the physical tasks distinguish the episode-specific requirement supplied by recollection from the familiar behavior supplied by the repertoire.

\subsubsection{Four Tasks and Success Criteria}
\label{app:realworld_tasks}

\paragraph{B2B and D2B: action reuse across layouts.}
The robot grasps a banana or duck from the table and places it inside the box.
Evaluation configurations vary object and box positions and the object's initial pose within the declared workspace.
Methods share 20 reset configurations in randomized order.
Success requires the correct object to be released fully inside the box within the episode limit, with a stable final placement and no corrective assistance.
These tasks test whether familiar repertoire actions remain useful as grasp and placement geometry changes; the simulation component comparisons quantify the contribution of access to the repertoire.

\paragraph{R2P: retain the original destination.}
A banana initially occupies one of three visually identical plates; the robot moves it to a common table-center staging position and then returns it to its original plate.
The instruction refers to the original plate without naming its location.
The 20-reset schedule allocates the three starting plates as evenly as possible (7, 7, and 6 trials), independently of layout and method, and randomizes trial order.
After staging, the three plates are empty and the banana occupies the same central region; successful completion requires both transfers in order and a stable release on the correct source plate.
Trials without a return attempt fail the full-task criterion.
Uniform random destination choice succeeds with probability $1/3$ before accounting for the two required transfers.

\paragraph{SOR: retain the order of completed interactions.}
The same banana and rubber duck used in B2B and D2B are placed one at a time into a common receiving area and then retrieved in the same order.
Placement follows trial-varying one-step cues; retrieval refers to the order actually completed.
The schedule uses each placement order in ten trials and assigns final object positions independently of order across the evaluation layouts.
We design the retrieval phase so that transient placement cues have left the current input and final object arrangement is independent of chronology.
Success requires both placements and both retrievals in the recorded order within the episode limit; the first order violation ends sequence success.
Logging requested and completed placements separately makes the reference order explicit when an earlier action fails.

\paragraph{Historical cues at the decision point.}
R2P's original plate and SOR's placement order are established before the corresponding return or retrieval decisions.
Staging removes the banana from its initial plate, and SOR's transient placement cues leave the current camera views before retrieval.
The robot replans from current observations while retaining each policy's native episode history; the task criteria above score the resulting destination and order decisions.

\subsubsection{Outcome Accounting and Reporting}
\label{app:realworld_results}

\begin{table}[!ht]
\centering\small
\caption{\textbf{Physical task structure and episode-specific decisions.}
Success counts retain the original 20 autonomous trials per method/task; ordered phases are part of the full-task criterion.
Task difficulty and the diversity of the training-only repertoire are distinct evaluation axes.}
\label{tab:realworld_dependency}
\begin{tabularx}{\linewidth}{@{}lXXc@{}}
\toprule
Task & Required manipulation sequence & Controlled variation or historical reference & WARM \\
\midrule
B2B / D2B & Pick up the named object, transport it, and release it stably inside the box. & Object/box layout and initial pose vary; objects are familiar. & 17/20; 16/20 \\
R2P & Source plate $\rightarrow$ common staging area $\rightarrow$ original source plate. & Three identical source plates occur 7/7/6 times; the return instruction does not name the plate. & 15/20 \\
SOR & Place two objects, then retrieve both in the completed placement order. & Both orders occur ten times; final object positions are assigned independently of chronology. & 15/20 \\
\bottomrule
\end{tabularx}
\end{table}

\paragraph{Task difficulty and retrieval ambiguity.}
R2P and SOR combine ordered manipulation with decisions whose reference is established earlier in the current episode.
The repertoire supplies familiar interaction structure, while recollection carries the destination or order needed for the current task.
These requirements make the physical study a test of history-dependent behavioral reuse, beyond successful grasping alone.
The repertoire contains 1,356 events from 20 training demonstrations, with five separate development demonstrations; the 20 autonomous test trials per task are counted independently.
A larger or more ambiguous repertoire tests a complementary property: whether consequence assessment remains useful when more retrieved experiences compete or the required skill is absent.
The present physical results establish the stated task and scene variations within the compact repertoire, without treating its event count as independent demonstration diversity.

Figure~\ref{fig:realworld}b reports autonomous successes out of 20 for each method and task under shared reset schedules.
Success follows the task-specific criteria above; timeouts, wrong-object actions, dropped objects, order violations, and corrective interventions count as failures.
The common denominator makes the four task comparisons directly interpretable, while the R2P destination and SOR order criteria distinguish history-dependent choices from physical execution.

\paragraph{Deployment sequences and retrieved cross-task events.}
Figure~\ref{fig:realworld_c4_sequences} extends the two examples in Figure~\ref{fig:realworld}a to seven physical conditions: B2B, D2B, R2P, and SOR on the original tabletop, plus D2B, R2P, and SOR on a black tabletop.
These sequences provide qualitative views of the four tasks; Figure~\ref{fig:realworld}b reports the 20-trial totals at task level, without separate per-appearance rates.
Each task row follows one deployment from the external/head view, with wrist-view insets; the adjacent \emph{Past recall} row shows observations indexed from retrieved events in other tasks.
The pairs connect cross-task interaction histories to the policy's physical execution under the current goal and scene.


\paragraph{Scene variation and repertoire reuse.}
On the original tabletop, B2B and D2B progress from object pickup to stable placement in the box.
The D2B sequence repeats that transfer on the black tabletop, where the background appearance and object--box arrangement change.
Across these examples, retrieved events pair banana and duck interactions with different destinations, making reusable grasp-and-transport structure visible across goals.
The deployments show \method{} carrying that familiar manipulation into the changed scenes while completing the current object-to-box goal.

\paragraph{Recollection in history-dependent tasks.}
In both R2P scenes, the banana leaves its initial plate, reaches the central staging area, and returns to the same plate.
In both SOR scenes, the robot completes the placements and retrieves the objects in their earlier placement order.
The object identities remain familiar across all four tasks, while the correct R2P destination and SOR retrieval order depend on the current episode's past.
These visible choices connect recollection's episode-specific requirements to the repertoire's familiar manipulation behavior in physical execution.

Together, the seven examples make the two roles of retrospection concrete: selected cross-task events supply reusable action structure, and recollection carries the destination or order established earlier in the current episode.
The 20-trial comparison strengthens this visual evidence: \method{} succeeds on 15/20 R2P and SOR trials, while each baseline reaches at most 9/20.

\subsection{Recorded-Observation Deployment Cost}
\label{app:recorded_cost}

We profile our shared-stage WARM checkpoint at update 3,000 and the released FastWAM checkpoint through their deployment implementations on one H100 80GB, with BF16 and batch size one.
Both receive the same 100 recorded current observations from five Press Button development episodes (20 replanning prefixes each), with 20 additional warmup queries per implementation.
We fix action horizon 32, execution prefix 4, and 10 flow evaluations.
Future observations are discarded before inference; WARM's episode memory receives the recorded executed actions, and generated actions are never executed.
This measures inference cost on a fixed observation cohort, rather than closed-loop success or a matched-training performance comparison. This additional profile is separate from the task-specialist accuracy protocol.

\begin{table}[h]
\centering\small
\caption{\textbf{Recorded-observation deployment cost.} CUDA-synchronized latency over the full replanning path after warmup. Memory is the maximum PyTorch allocated/reserved GPU memory over 100 queries, excluding CPU bank storage and system RAM.}
\label{tab:recorded_cost}
\begin{tabular}{@{}lrrrr@{}}
\toprule
Implementation & Median (ms) & p95 (ms) & Allocated (GiB) & Reserved (GiB) \\
\midrule
FastWAM release & 366.3 & 381.8 & 23.31 & 23.58 \\
WARM & 530.1 & 565.2 & 23.62 & 23.74 \\
\bottomrule
\end{tabular}
\end{table}

The WARM bank contains 89,680 events, with exact CPU cosine retrieval and $K=32$; gathered candidates are transferred to the model device.
Recent-event and action-summary capacities are 20 and 8, respectively, while their observed occupancies are 0--8 and 0--8.
The recent-event buffer does not saturate in this cohort, so this is not a worst-case history-length measurement.
The median latency ratio is 1.447 (WARM/FastWAM).


\paragraph{Nested retrieval intervals.}
We additionally reconcile the retained retrieval stopwatches with all 100 synchronized outer measurements, with identical recorded sample keys and the declared warmup excluded.
Table~\ref{tab:retrieval_phase_cost} reports these inclusive intervals.
The per-query outer time minus the complete retrieval interval defines the remaining replanning work; it includes observation processing, causal history, source construction, sampling, and telemetry, and does not isolate WARM's adaptation, consequence, or gate modules.
Only means are additive across this nested retrieval/remainder partition.

\begin{table}[htbp]
\centering\small
\caption{\textbf{Nested retrieval stopwatch accounting.} The same 100 recorded DEV prefixes and shared-stage update-3,000 checkpoint as Table~\ref{tab:recorded_cost}. Retrieval subintervals are included in complete retrieval, which is included in full replanning. These inclusive intervals describe the measured computation.}
\label{tab:retrieval_phase_cost}
\begin{tabular}{@{}lrrr@{}}
\toprule
Stopwatch interval & Mean (ms) & Median (ms) & p95 (ms) \\
\midrule
Full replanning & 524.27 & 530.06 & 565.22 \\
Complete retrieval & 39.11 & 38.44 & 43.65 \\
\quad Image preprocessing & 9.36 & 9.21 & 10.54 \\
\quad Frozen DINO encoding & 6.03 & 5.95 & 6.61 \\
\quad Exact CPU bank search & 22.09 & 21.34 & 25.96 \\
\quad Candidate gather & 0.14 & 0.13 & 0.15 \\
\quad Remaining retrieval work & 1.49 & 1.49 & 1.53 \\
All other replanning work & 485.16 & 488.10 & 526.19 \\
\bottomrule
\end{tabular}
\end{table}

\paragraph{Timing boundary and initialization.}
We synchronize CUDA immediately before and after the whole deployed replanning call, including image/state preprocessing, retrieval and memory operations where present, generation, and queue publication.
Disk decoding and simulator execution are excluded.
Both implementations use a $[1,3,384,320]$ image composite with their native numerical preprocessing.
WARM's call includes built-in policy telemetry; external evidence writing for the FastWAM reference is outside its timed call.
The latency difference therefore compares the deployed implementations, including this instrumentation difference, and is not a module-level attribution.
One-time policy construction, including strict artifact checks and loading on the AFS-backed deployment, takes 854.9\,s for WARM and 144.3\,s for FastWAM in this run; it is recorded separately from query latency.
No other GPU compute process is observed in either run at five-second sampling; this observation is not a hardware-exclusivity guarantee.

\paragraph{Scope.}
The FastWAM row is a release-checkpoint compute reference, not a same-data fine-tuned accuracy baseline.
This profile does not establish a success--cost trade-off, MemoryWAM's cost, robot control frequency, or the cost with a fully occupied 20-event recent memory.
Those comparisons require their own matched experiments.




\subsection{Matched-Model Recorded-Observation Cost}
\label{app:matched_cost}

We measure independently trained Full and Pooled models at fixed final update 3,000 (seed 3407), using the same base checkpoint, data sources, optimizer/scheduler and update budget.
Pooled learns four memory tokens from all $K$ retrieved action, visual/effect, timing and context records alongside the same eight causal-history gist tokens.
It removes explicit adaptation, consequence scoring and influence gating; the four associated retrieval/effect/adaptation/gate losses are zero, and other objectives remain.
This compares complete methods with their supervision, not an isolated same-objective architecture ablation or a named-baseline reproduction.

Both receive the same 100 observations from five Press Button development episodes (20 prefixes each), with 20 additional warmups per method.
We use one H100 80GB, BF16, batch one, $H=32$, execution prefix four, NFE ten, $K=32$, and recent/summary capacities 20/8.
Future observations are discarded; only recorded executed actions update history, and generated commands are never executed. The bank has 89,680 events.
CUDA synchronization encloses the complete deployed replanning call, including preprocessing, retrieval, history, generation, telemetry and queue publication; disk decoding, model construction and simulation are excluded.

\begin{table}[htbp]
\centering\small
\caption{\textbf{Matched-model full-replanning cost.} The same fixed observation schedule and budget, with separate warmups and prespecified Full-then-Pooled execution. Allocated/reserved memory is the measured peak PyTorch GPU memory, excluding CPU bank storage and system RAM.}
\label{tab:matched_recorded_cost}
\begin{tabular}{@{}lrrrr@{}}
\toprule
Method & Median (ms) & p95 (ms) & Allocated (GiB) & Reserved (GiB) \\
\midrule
Full & 473.4 & 500.3 & 23.62 & 23.74 \\
Pooled & 476.5 & 501.1 & 23.61 & 23.75 \\
\bottomrule
\end{tabular}
\end{table}

The Full/Pooled median ratio is 0.993: costs are close on this cohort, without establishing a statistical speedup, worst-case history cost or causal per-module attribution.
One-time construction, including strict validation and AFS-backed loading, takes 439.0/420.4\,s for Full/Pooled and is excluded from query timings.
These models and measurements are separate from the earlier WARM/released-FastWAM reference in Table~\ref{tab:recorded_cost}. They measure computation, not task success, and do not replace the submitted task-specialist accuracy protocol.


\paragraph{Instrumented component intervals.}
We separately observe the same Full model on these 100 prefixes after 20 warmups, using passive forward hooks and method observers.
Each of the five regions in Table~\ref{tab:matched_module_intervals} is invoked once in all 100 measured queries.
We record host wall and CUDA-stream event intervals; event resolution and evidence writing follow the outer query timings. Observers preserve method returns, parameters and the sampler.

\begin{table}[htbp]
\centering\small
\caption{\textbf{Observed Full component intervals (ms).} Medians and p95 over 100 queries. Event adaptation includes effect prediction and action-context construction; source resolution includes these regions, candidate ranking, the confidence gate and source construction. The confidence-gate row excludes subsequent eligibility/blending operations. Parent/child intervals overlap; observer enqueue overhead and host gaps may contribute. These are not exclusive kernel sums or causal removal costs.}
\label{tab:matched_module_intervals}
\begin{tabularx}{\linewidth}{@{}Xrrrr@{}}
\toprule
Region & Host median & Host p95 & CUDA median & CUDA p95 \\
\midrule
Source resolution (inclusive) & 14.940 & 16.166 & 14.895 & 16.107 \\
Event adapter (inclusive) & 3.170 & 3.327 & 3.133 & 3.290 \\
Effect prediction & 0.600 & 0.644 & 0.565 & 0.610 \\
Candidate reranker & 0.548 & 0.579 & 0.530 & 0.561 \\
Confidence-gate forward & 0.723 & 0.771 & 0.692 & 0.739 \\
\bottomrule
\end{tabularx}
\end{table}

The instrumented outer median/p95 is 477.2/506.5\,ms, reported separately from Table~\ref{tab:matched_recorded_cost}; the runs do not isolate observer overhead.
Nested intervals or quantiles must not be summed. These costs do not establish task-success contributions or natural skill-absent reliability.
